% !TEX TS-program = XeLaTeX
% use the following command:
% all document files must be coded in UTF-8
\documentclass[english]{textolivre}
% build HTML with: make4ht -e build.lua -c textolivre.cfg -x -u article "fn-in,svg,pic-align"

\journalname{Texto Livre}
\thevolume{19}
%\thenumber{1} % old template
\theyear{2026}
\receiveddate{\DTMdisplaydate{2025}{9}{9}{-1}} % YYYY MM DD
\accepteddate{\DTMdisplaydate{2025}{10}{24}{-1}}
\publisheddate{\DTMdisplaydate{2026}{4}{13}{-1}}
\corrauthor{Jamil Suprihatiningrum}
\articledoi{10.1590/1983-3652.2026.61658}
%\articleid{NNNN} % if the article ID is not the last 5 numbers of its DOI, provide it using \articleid{} commmand 
% list of available sesscions in the journal: articles, dossier, reports, essays, reviews, interviews, editorial
\articlesessionname{articles}
\runningauthor{Suprihatiningrum and Syaripudin} 
%\editorname{Leonardo Araújo} % old template
\sectioneditorname{Daniervelin Pereira~\orcid{0000-0003-1861-3609}}
\layouteditorname{Saula Cecília~\orcid{0009-0006-3069-8480}}

\title{Technology for equity: a systematic review of AI, adaptive tools, and digital platforms in differentiated instruction}
\othertitle{Tecnologia para a equidade: uma revisão sistemática de IA, ferramentas adaptativas e plataformas digitais na instrução diferenciada}
% if there is a third language title, add here:
%\othertitle{Artikelvorlage zur Einreichung beim Texto Livre Journal}

\author[1]{Jamil Suprihatiningrum~\orcid{0000-0002-1836-8845}\thanks{Email: \href{mailto:jamil.suprihatiningrum@uin-suka.ac.id}{jamil.suprihatiningrum@uin-suka.ac.id}}}
\author[1]{Ahmad Syaripudin~\orcid{0009-0000-5606-0406}\thanks{Email: \href{mailto:24204081005@student.uin-suka.ac.id}{24204081005@student.uin-suka.ac.id}}}
\affil[1]{State Islamic University Sunan Kalijaga, Yogyakarta, Indonesia.}

\addbibresource{article.bib}
% use biber instead of bibtex
% $ biber article

% used to create dummy text for the template file
\definecolor{dark-gray}{gray}{0.35} % color used to display dummy texts
\usepackage{lipsum}
\SetLipsumParListSurrounders{\colorlet{oldcolor}{.}\color{dark-gray}}{\color{oldcolor}}

% used here only to provide the XeLaTeX and BibTeX logos
\usepackage{hologo}

% if you use multirows in a table, include the multirow package
\usepackage{multirow}

% provides sidewaysfigure environment
\usepackage{rotating}

% CUSTOM EPIGRAPH - BEGIN 
%%% https://tex.stackexchange.com/questions/193178/specific-epigraph-style
\usepackage{epigraph}
\renewcommand\textflush{flushright}
\makeatletter
\newlength\epitextskip
\pretocmd{\@epitext}{\em}{}{}
\apptocmd{\@epitext}{\em}{}{}
\patchcmd{\epigraph}{\@epitext{#1}\\}{\@epitext{#1}\\[\epitextskip]}{}{}
\makeatother
\setlength\epigraphrule{0pt}
\setlength\epitextskip{0.5ex}
\setlength\epigraphwidth{.7\textwidth}
% CUSTOM EPIGRAPH - END

% to use IPA symbols in unicode add
%\usepackage{fontspec}
%\newfontfamily\ipafont{CMU Serif}
%\newcommand{\ipa}[1]{{\ipafont #1}}
% and in the text you may use the \ipa{...} command passing the symbols in unicode

% LANGUAGE - BEGIN
% ARABIC
% for languages that use special fonts, you must provide the typeface that will be used
% \setotherlanguage{arabic}
% \newfontfamily\arabicfont[Script=Arabic]{Amiri}
% \newfontfamily\arabicfontsf[Script=Arabic]{Amiri}
% \newfontfamily\arabicfonttt[Script=Arabic]{Amiri}
%
% in the article, to add arabic text use: \textlang{arabic}{ ... }
%
% RUSSIAN
% for russian text we also need to define fonts with support for Cyrillic script
% \usepackage{fontspec}
% \setotherlanguage{russian}
% \newfontfamily\cyrillicfont{Times New Roman}
% \newfontfamily\cyrillicfontsf{Times New Roman}[Script=Cyrillic]
% \newfontfamily\cyrillicfonttt{Times New Roman}[Script=Cyrillic]
%
% in the text use \begin{russian} ... \end{russian}
% LANGUAGE - END

% EMOJIS - BEGIN
% to use emoticons in your manuscript
% https://stackoverflow.com/questions/190145/how-to-insert-emoticons-in-latex/57076064
% using font Symbola, which has full support
% the font may be downloaded at:
% https://dn-works.com/ufas/
% add to preamble:
% \newfontfamily\Symbola{Symbola}
% in the text use:
% {\Symbola }
% EMOJIS - END

% LABEL REFERENCE TO DESCRIPTIVE LIST - BEGIN
% reference itens in a descriptive list using their labels instead of numbers
% insert the code below in the preambule:
%\makeatletter
%\let\orgdescriptionlabel\descriptionlabel
%\renewcommand*{\descriptionlabel}[1]{%
%  \let\orglabel\label
%  \let\label\@gobble
%  \phantomsection
%  \edef\@currentlabel{#1\unskip}%
%  \let\label\orglabel
%  \orgdescriptionlabel{#1}%
%}
%\makeatother
%
% in your document, use as illustraded here:
%\begin{description}
%  \item[first\label{itm1}] this is only an example;
%  % ...  add more items
%\end{description}
% LABEL REFERENCE TO DESCRIPTIVE LIST - END


% add line numbers for submission
%\usepackage{lineno}
%\linenumbers

\begin{document}
\maketitle

\begin{polyabstract}
\begin{abstract}
This study examines how artificial intelligence (AI), adaptive learning systems, and digital tools enhance differentiated instruction (DI) to advance educational equity in digital classrooms. Using a systematic review of 41 peer-reviewed studies published between 2015 and 2025, the study analyzes how different technologies support the four components of DI, namely content, process, product, and learning environment, under the PRISMA protocol. The synthesis identifies four main themes. Digital tools increase accessibility and foster multimodal engagement, though their adaptive functions often remain static. AI and adaptive systems enable real-time personalization but face persistent challenges related to algorithmic transparency, data demands, and contextual relevance. Innovative products and immersive environments such as virtual reality labs, educational games, and multimodal storytelling extend inclusion by providing multiple response formats, text-to-speech, and flexible submission options aligned with Universal Design for Learning principles. Teacher attitudes and institutional support consistently determine whether technological innovations lead to sustainable differentiation. The review demonstrates that technology strengthens DI only when educators integrate it with professional judgment, institutional commitment, and systemic equity frameworks. The findings offer actionable insights for teachers, researchers, and policymakers seeking to design digital learning environments that promote inclusion and educational justice.

\keywords{Differentiated instruction\sep Artificial Intelligence\sep Adaptive learning systems\sep Digital platforms\sep Equity in education}
\end{abstract}

\begin{portuguese}
\begin{abstract}
Este estudo examina como a inteligência artificial (IA), os sistemas de aprendizagem adaptativa e as ferramentas digitais fortalecem a instrução diferenciada (ID) para promover a equidade educacional em salas de aula digitais. Com base em uma revisão sistemática de 41 estudos revisados por pares publicados entre 2015 e 2025, o estudo analisa como diferentes tecnologias apoiam os quatro componentes da ID, a saber, conteúdo, processo, produto e ambiente de aprendizagem, seguindo o modelo PRISMA. A síntese identifica quatro temas principais. As ferramentas digitais aumentam a acessibilidade e promovem o engajamento multimodal, embora suas funções adaptativas frequentemente permaneçam estáticas. A IA e os sistemas adaptativos permitem a personalização em tempo real, mas enfrentam desafios persistentes relacionados à transparência algorítmica, à demanda de dados e à relevância contextual. Produtos inovadores e ambientes imersivos, como laboratórios de realidade virtual, jogos educacionais e narrativas multimodais, ampliam a inclusão ao oferecer múltiplos formatos de resposta, conversão de texto em fala e opções flexíveis de envio, de acordo com os princípios do Desenho Universal para a Aprendizagem. As atitudes dos professores e o apoio institucional determinam de forma consistente se as inovações tecnológicas levam à diferenciação sustentável. No geral, a revisão demonstra que a tecnologia fortalece a ID somente quando os educadores a integram ao julgamento profissional, ao compromisso institucional e aos marcos sistêmicos de equidade. Os resultados oferecem orientações práticas para professores, pesquisadores e formuladores de políticas que buscam projetar ambientes de aprendizagem digitais que promovam a inclusão e a justiça educacional.


\keywords{Instrução diferenciada\sep Inteligência Artificial\sep Sistemas de aprendizagem adaptativa\sep Plataformas digitais\sep Equidade na educação}
\end{abstract}
\end{portuguese}
% if there is another abstract, insert it here using the same scheme
\end{polyabstract}

\section{Introduction}\label{sec-intro}
Differentiated instruction (DI) is now widely acknowledged as crucial for meeting learner diversity. It addresses differences in readiness, interests, learning preferences, and sociocultural backgrounds by adjusting content, process, product, and the learning environment \cite{tomlinson2001, tomlinson2014}. Instead of a one-size-fits-all model, DI emphasizes flexibility in both design and implementation to foster equity, engagement, and better learning outcomes \cite{ferrer2025, puzio2020, ramaila2025, sabir2024}. Globally, including Indonesia, DI has been officially supported through curriculum reforms and inclusive policies aimed at reaching all learners, especially those from marginalized or linguistically diverse backgrounds \cite{marlina2019, ouyang2023, smalejacobse2019}. Its value lies in its ability to reduce opportunity gaps, promote learner independence, and develop positive attitudes toward learning \cite{cullen2024, donald2022, saad2025}. However, implementing DI can be difficult and requires teacher expertise, careful planning, and effective classroom management \cite{maulana2019, vangeel2019}. Many teachers struggle to maintain individualized learning paths, track student progress, and manage diverse needs in large classrooms \cite{cevikbas2022, darlinghammond2020,kariippanon2019, kerimbayev2025, tetzlaff2020}. These challenges have increased interest in using technology to make DI more scalable, efficient, and adaptable across different settings.

Over the past two decades, digital transformation has facilitated a gradual shift toward more personalized learning models \cite{engelbrecht2020, menaguacas2025, shemshack2020, tetzlaff2020}. The rise of artificial intelligence (AI), adaptive learning platforms, and interactive tools has created new opportunities for responsive instruction. AI systems, including intelligent tutors and recommender engines, evaluate student performance in real time and adjust content dynamically \cite{gligorea2023, naseer2024, sajja2024, sayed2022, taskin2025}. Adaptive platforms organize tasks based on learner input, enabling individualized pacing, while digital tools like learning management systems, educational games, and collaborative applications create multimodal environments that accommodate diverse learning preferences \cite{jose2024, strielkowski2024, vincentruz2022, wang2020}. These technologies reduce the practical challenges teachers face when implementing DI because they automate assessments, provide tailored pathways, and enhance student engagement \cite{krishan2023, palieraki2021, vaganova2020}. When aligned with DI principles, they have the potential to support large-scale personalization, even though issues related to contextual fit, ethics, and equity remain unresolved.

Theoretically, DI and educational technologies share constructivist, learner-centered foundations emphasizing scaffolding, mastery, and self-regulation \cite{damyanov2024, donkoh2024, guo2024}. Adaptive platforms reflect Vygotsky’s concept of the zone of proximal development by offering targeted support \cite{damyanov2024, donkoh2024, guo2024}, while AI-driven personalization functions as formative feedback that sequences tasks based on learner performance. In practice, digital tools can enact DI principles by enabling flexible pacing, varied content, and real-time monitoring. However, the level of alignment depends on how these tools are designed and implemented. Some applications genuinely extend the logic of DI, while others offer superficial personalization with fixed content or opaque algorithms \cite{burner2025, chang2023, kochmar2021, liao2024, murtaza2022}. For technology and DI to work together effectively, their integration must be based on pedagogical frameworks prioritizing teacher agency and equitable outcomes.

Empirical evidence on technology-supported DI remains mixed, varying in strength, setting, and technological scope. Some studies show positive effects on formative assessment and motivation, while others indicate limited use or specific challenges in certain contexts \cite{burner2025, sree2024}. Adaptive systems might benefit lower-achieving students more than high achievers, and dashboards heavily depend on teacher interpretation \cite{molenaar2019, pozdniakov2025}. Most research focuses on short-term outcomes, leaving long-term factors like self-regulation, motivation, and inclusion less examined \cite{farrokhnia2025, held2024, zafar2024}. These differences highlight the need for a thorough review that categorizes technologies, clarifies the conditions that support them, and situates DI within broader equity frameworks.

Previous reviews have examined aspects of DI and technology separately, often emphasizing digital personalization or AI ethics without exploring their intersection. For instance, \textcite{smalejacobse2019} focused on DI implementation in traditional classrooms, while \textcite{shemshack2020} reviewed personalized learning terminology. \textcite{engelbrecht2020} examined mathematics classrooms during digital transformation, and \textcite{damyanov2024} addressed special education. In Latin America, some reviews have explored AI and digital technologies, but with limited pedagogical focus. \textcite{salaspilco2022} mapped AI applications in higher education, identifying adaptive learning trends but overlooking how they align with DI. \textcite{silvaleon2024} analyzed broader digital innovations, emphasizing entrepreneurship and social inclusion rather than classroom differentiation. \textcite{portilla2025} provided a global synthesis of AI in Education 4.0, highlighting algorithmic bias and personalization challenges without situating them in equity-oriented pedagogies. Meanwhile, \textcite{carrillo2026} explored AI’s structural and policy implications in Latin American higher education but did not address adaptive or differentiated learning processes. None of these reviews systematically investigates how AI, adaptive tools, and digital platforms jointly operationalize DI to promote equity.

Building on this gap, the present review examines how these technologies are integrated into DI and assesses how such integration advances educational equity across its four components (content, process, product, and environment). To achieve this aim, the review addresses three guiding research questions.

\begin{enumerate}
    \item RQ1. How are AI, adaptive systems, and digital platforms used to implement DI across educational contexts?
    \item RQ2. Which pedagogical, institutional, and contextual factors enable or constrain equitable technology-supported DI?
    \item RQ3. In what ways do these technologies contribute to advancing educational equity through DI?
\end{enumerate}

\section{Method}

\subsection{Research design}
This study followed a qualitative Systematic Literature Review using the PRISMA 2020 framework \cite{page2021}, as seen in Figure \ref{fig-1}, to ensure transparency and reproducibility. It synthesized empirical studies from 2015 to 2025 on integrating AI, adaptive learning systems, and digital tools within DI. This period reflects a decade of rapid growth in AI and adaptive technologies in education.

The scope of ``AI'' in this review includes intelligent tutoring systems, machine learning–based recommender engines, and generative AI applications. ``Adaptive systems'' refer to algorithmic or data-driven platforms that personalize learning processes based on learner input or performance data. ``Digital tools'' encompass multimodal and collaborative technologies used in classroom contexts, such as virtual labs, educational games, and online learning platforms.

%--- codigo da figura 1 ---%
\begin{figure}[h!]
\centering
\begin{minipage}{.90\textwidth}
\includegraphics[width=\textwidth]{images/Figure1.png}
\caption{The review protocol.}
\label{fig-1}
\source{Own elaboration.}
\end{minipage}
\end{figure}

\subsection{Data sources and search strategy}
A comprehensive search was performed across six databases: Scopus, Web of Science, ERIC, ProQuest, ScienceDirect, and SpringerLink. The search strategy combined keywords and Boolean operators to identify studies on DI integrated with AI, adaptive learning systems, and digital technologies. The final search string used in all databases was:

\begin{quote}
    (``differentiated instruction" OR ``differentiated teaching") AND (``artificial intelligence" OR ``AI" OR ``adaptive learning" OR ``adaptive system*" OR ``digital platform*" OR ``digital tool*" OR ``educational technolog*" OR ``technology integration") AND (``kindergarten" OR ``early childhood" OR ``elementary" OR ``primary" OR ``secondary" OR ``K-12" OR ``higher education" OR ``university" OR ``undergraduate*" OR ``vocational training").
\end{quote}

The formula used the wildcard symbol (*) to include word variations and linked conceptual groups with AND operators to ensure inclusive retrieval of studies on AI. Searches were limited to the title, abstract, and keyword fields and included educational level terms (e.g., kindergarten, K-12, higher education). All records were managed in EndNote, duplicates removed, and results verified through backward and forward citation tracking before applying the inclusion and exclusion criteria.

\subsection{Inclusion and exclusion criteria}
The review applied specific criteria to ensure rigor and relevance, summarized in Table \ref{tab-1}. To maintain consistent data extraction and accurate interpretation of findings, only studies published in English were included. This decision was made because English dominates peer-reviewed academic publishing, and reliable translations are limited across databases. Although this may introduce language bias, it helped keep coding and quality assessment procedures consistent across all sources. The possible effects of this limitation are discussed in the final section of the paper.

%--- codigo da tabela 1 ---%
\begin{table}[htbp]
\centering
\begin{threeparttable}
\caption{Inclusion and Exclusion Criteria.}\label{tab-1}
\begin{tabular}{p{2.7cm}p{5.5cm}p{5.2cm}}
\toprule
Criterion & Inclusion & Exclusion \\
\midrule
Publication type & Peer-reviewed empirical research (qualitative, quantitative, mixed methods) & Non-peer-reviewed reports, theoretical papers, or conference abstracts without full papers, grey literature \\
Publication period & 2015--2025 & Before 2015 or after 2025 \\
Language & English & Non-English publications \\
Focus & Studies addressing DI combined with AI, adaptive systems, or digital platforms & Studies unrelated to DI or limited to conventional/non-digital instruction \\
Educational level & Early childhood, primary, secondary, higher, or vocational education & Studies not specifying or irrelevant to educational contexts \\
Availability & Full-text accessible & Abstract-only or inaccessible full text \\
\bottomrule
\end{tabular}
\source{Own elaboration.}
\end{threeparttable}
\end{table}

\subsection{Screening and selection process}
After removing duplicates, two reviewers independently screened the titles and abstracts against the inclusion criteria. Potentially eligible studies then underwent full-text review, and any disagreements were resolved through discussion to ensure consistency. The process and exclusion decisions at the full-text stage were documented to maintain transparency.

\subsection{Data extraction and analysis}
Data were systematically extracted using a structured template that covered: bibliographic details; study context; technology type; DI components; research design; and key findings. The coding system for thematic analysis was developed inductively and refined through iterative comparison. Initial codes based on Tomlinson’s \citeyear{tomlinson1999} four DI dimensions were aligned with technology categories, pilot-tested on ten studies, and collaboratively refined into a final codebook. Thematic coding was supported by an AI-assisted text analysis tool to organize and cluster data during the initial synthesis stage. All coding decisions and thematic interpretations were manually verified by the researchers to ensure analytical consistency and validity, with each step documented for replicability and auditability.

We also constructed a co-occurrence matrix from the final codes (technology category, DI component, mediators, and equity dimensions) and created a thematic network (as presented in Figure \ref{fig-4}) to visualize the connections among topics. An AI-assisted tool helped with early clustering, and all links were checked manually.

\section{Findings and discussion}
This section presents the main findings and discusses their implications. To give an overview, Table \ref{tab-2} maps the selected studies based on the types of technology used (AI-based systems, adaptive tools, and broader digital platforms) and the DI components they cover (content, process, product, and environment).

%--- codigo da tabela 2 ---%
\begin{table}[h!]
\centering
\small
\caption{Studies mapped by technology type and DI component.}\label{tab-2}
\resizebox{\textwidth}{!}{
\begin{tabular}{c p{5cm} c c c c c c c}
\toprule
No. & Authors & AI & Adaptive Tools & Digital Tool/Platform & Content & Process & Product & Environment \\
\midrule
1 & \textcite{meutstege2023} & - & $\sqrt{}$ & $\sqrt{}$ & $\sqrt{}$ & $\sqrt{}$ & $\sqrt{}$ & $\sqrt{}$ \\
2 & \textcite{hasanah2022} & - & $\sqrt{}$ & $\sqrt{}$ & $\sqrt{}$ & $\sqrt{}$ & $\sqrt{}$ & - \\
3 & \textcite{min2018} & - & $\sqrt{}$ & $\sqrt{}$ & $\sqrt{}$ & $\sqrt{}$ & $\sqrt{}$ &  \\
4 & \textcite{valiandes2018} & - & - & $\sqrt{}$ & $\sqrt{}$ & $\sqrt{}$ & $\sqrt{}$ & $\sqrt{}$ \\
5 & \textcite{smets2020} & - & - & $\sqrt{}$ & $\sqrt{}$ & $\sqrt{}$ & $\sqrt{}$ & $\sqrt{}$ \\
6 & \textcite{westbroek2020} & - & - & $\sqrt{}$ & $\sqrt{}$ & $\sqrt{}$ & $\sqrt{}$ & $\sqrt{}$ \\
7 & \textcite{strogilos2021} & - & $\sqrt{}$ & $\sqrt{}$ & $\sqrt{}$ & $\sqrt{}$ & $\sqrt{}$ & - \\
8 & \textcite{grifulfreixenet2021} & - & - & $\sqrt{}$ & $\sqrt{}$ & $\sqrt{}$ & $\sqrt{}$ & - \\
9 & \textcite{vantieghem2020} & - & - & $\sqrt{}$ & $\sqrt{}$ & $\sqrt{}$ & $\sqrt{}$ & - \\
10 & \textcite{valiandes2015} & - & $\sqrt{}$ & - & $\sqrt{}$ & $\sqrt{}$ & $\sqrt{}$ & - \\
11 & \textcite{karst2022} & - & $\sqrt{}$ & $\sqrt{}$ & $\sqrt{}$ & $\sqrt{}$ & - & - \\
12 & \textcite{forster2018} & - & $\sqrt{}$ & $\sqrt{}$ & $\sqrt{}$ & $\sqrt{}$ & - & - \\
13 & \textcite{iterbeke2020} & - & - & $\sqrt{}$ & $\sqrt{}$ & $\sqrt{}$ & - & - \\
14 & \textcite{maulana2019} & - & - & $\sqrt{}$ & $\sqrt{}$ & $\sqrt{}$ & $\sqrt{}$ & - \\
15 & \textcite{songer2019} & $\sqrt{}$ & $\sqrt{}$ & $\sqrt{}$ & $\sqrt{}$ & $\sqrt{}$ & $\sqrt{}$ & - \\
16 & \textcite{wu2019} & - & $\sqrt{}$ & $\sqrt{}$ & $\sqrt{}$ & $\sqrt{}$ & $\sqrt{}$ & - \\
17 & \textcite{easa2023} & - & - & $\sqrt{}$ & $\sqrt{}$ & $\sqrt{}$ & $\sqrt{}$ & - \\
18 & \textcite{hasanah2023} & - & - & $\sqrt{}$ & $\sqrt{}$ & $\sqrt{}$ & $\sqrt{}$ & - \\
19 & \textcite{halil2024} & - & - & $\sqrt{}$ & $\sqrt{}$ & $\sqrt{}$ & $\sqrt{}$ & - \\
20 & \textcite{zervoudakis2019} & $\sqrt{}$ & - & $\sqrt{}$ & - & $\sqrt{}$ & - & - \\
21 & \textcite{cihad2018} & - & - & $\sqrt{}$ & $\sqrt{}$ & $\sqrt{}$ & $\sqrt{}$ & - \\
22 & \textcite{estaiteyeh2023} & - & - & $\sqrt{}$ & $\sqrt{}$ & $\sqrt{}$ & $\sqrt{}$ & - \\
23 & \textcite{spyropoulou2025} & - &  & $\sqrt{}$ & $\sqrt{}$ & $\sqrt{}$ & $\sqrt{}$ & $\sqrt{}$ \\
24 & \textcite{shareefa2024} & - & - & $\sqrt{}$ & $\sqrt{}$ & $\sqrt{}$ & $\sqrt{}$ & $\sqrt{}$ \\
25 & \textcite{shareefa2023} &  &  &  &  &  &  &  \\
26 & \textcite{yuen2023} & - & - & $\sqrt{}$ & $\sqrt{}$ & $\sqrt{}$ & $\sqrt{}$ & $\sqrt{}$ \\
27 & \textcite{malacapay2019} & - & - & $\sqrt{}$ & $\sqrt{}$ & $\sqrt{}$ & $\sqrt{}$ & - \\
28 & \textcite{gibbs2021} & - & - & $\sqrt{}$ & $\sqrt{}$ & $\sqrt{}$ & $\sqrt{}$ & $\sqrt{}$ \\
29 & \textcite{heng2023} & - & $\sqrt{}$ & - & $\sqrt{}$ & $\sqrt{}$ & $\sqrt{}$ & $\sqrt{}$ \\
30 & \textcite{smets2018} & - & - & $\sqrt{}$ & $\sqrt{}$ & $\sqrt{}$ & $\sqrt{}$ & $\sqrt{}$ \\
31 & \textcite{vaganova2020} & - & $\sqrt{}$ & $\sqrt{}$ & $\sqrt{}$ & $\sqrt{}$ & $\sqrt{}$ & - \\
32 & \textcite{palieraki2021} & - & $\sqrt{}$ & $\sqrt{}$ & $\sqrt{}$ & $\sqrt{}$ & $\sqrt{}$ & $\sqrt{}$ \\
33 & \textcite{naseer2024} & $\sqrt{}$ & $\sqrt{}$ &  & $\sqrt{}$ & $\sqrt{}$ & $\sqrt{}$ & - \\
34 & \textcite{krishan2023} & - & - & $\sqrt{}$ & $\sqrt{}$ & $\sqrt{}$ & $\sqrt{}$ & - \\
35 & \textcite{howorth2024} & $\sqrt{}$ & $\sqrt{}$ & $\sqrt{}$ & $\sqrt{}$ & $\sqrt{}$ & $\sqrt{}$ & - \\
36 & \textcite{kong2024} & $\sqrt{}$ & - & - & $\sqrt{}$ & $\sqrt{}$ & $\sqrt{}$ & - \\
37 & \textcite{kempe2019} & - & $\sqrt{}$ & $\sqrt{}$ & $\sqrt{}$ & $\sqrt{}$ & $\sqrt{}$ & - \\
38 & \textcite{yunianto2024} & $\sqrt{}$ & - & - & $\sqrt{}$ & $\sqrt{}$ & $\sqrt{}$ & - \\
39 & \textcite{woo2024} & $\sqrt{}$ & - & - & $\sqrt{}$ & $\sqrt{}$ & $\sqrt{}$ & - \\
40 & \textcite{chen2025} & $\sqrt{}$ & $\sqrt{}$ & $\sqrt{}$ & $\sqrt{}$ & $\sqrt{}$ & $\sqrt{}$ & - \\
41 & \textcite{sajja2024} & $\sqrt{}$ & $\sqrt{}$ & $\sqrt{}$ & $\sqrt{}$ & $\sqrt{}$ & $\sqrt{}$ & - \\
\bottomrule
\end{tabular}
}
\source{Own elaboration.}
\end{table}

Across the reviewed period, the early studies (2015-2019) primarily focused on digital tools for differentiated content and formative assessment, reflecting the initial stages of classroom digitalization. In contrast, studies from 2020-2025 increasingly incorporated AI and adaptive learning systems, emphasizing personalization, data analytics, and equity considerations. This shift suggests a gradual evolution from tool-based differentiation toward intelligent, context-sensitive systems that align with equity and inclusion frameworks. These patterns help clarify how tools fit within DI frameworks and guide the synthesis below.

\subsection{Publication and context overview}
Research on technology-enhanced DI increased significantly across 41 studies after 2018 (Figure \ref{fig-2}), aligning with digital transformation and remote learning. By level (Figure \ref{fig-3}), the collection includes eight primary, 14 secondary, five higher/vocational, five teacher-education, and nine cross-level K-12 investigations. Geographically, the studies cover 23 countries across Europe, Asia, North America, and Oceania, allowing for cautious generalization while remaining aware of policy and infrastructure factors that influence implementation.

%---codigo da figuras 2 e 3 ---%
\begin{figure}[h!]
\centering
\begin{minipage}{.80\textwidth}
\includegraphics[width=\textwidth]{images/Figure2.png}
\caption{Number of publications by year.}
\label{fig-2}
\source{Own elaboration.}
\end{minipage}
\end{figure}

\begin{figure}[h!]
\centering
\begin{minipage}{.80\textwidth}
\includegraphics[width=\textwidth]{images/Figure3.png}
\caption{Number of publications by level of education.}
\label{fig-3}
\source{Own elaboration.}
\end{minipage}
\end{figure}

\subsection{Technology types used in DI}
The review identifies four broad categories of technologies:

\begin{enumerate}
\item AI and machine learning (8 studies), including generative AI for writing support and feedback, intelligent assistants, optimization algorithms for grouping, and deep learning for adaptive pathways;
\item Adaptive learning platforms (8 studies) such as Learning Navigator, computer-based Learning Progress Assessment (LPA), and assessment-linked systems (e.g., Lernstand 5) that tailor instruction to readiness and performance;
\item General digital platforms and tools (20 studies), including digital games, Virtual Reality (VR) labs, ICT labs, online collaboration tools (Padlet, Google Classroom, Quizizz), and multimedia/storytelling applications;
\item Conceptual/minimal-integration discussions (7 studies) focusing on DI frameworks rather than concrete tools.
\end{enumerate}

\subsection{Network visualization of technology-DI relationships}
We generated a co-occurrence network in VOSviewer (v1.6.20) using the coded relationships among technology categories (AI, adaptive systems, and general digital tools), DI components (content, process, product, and environment), mediating factors (e.g., teacher mindset, professional development, infrastructure access, classroom constraints), and equity dimensions (e.g., accessibility, personalization, agency/autonomy). The network (Figure \ref{fig-4}) consists of 17 items connected through 136 weighted links. In the map, node size reflects how frequently each item appears across the 41 studies, edge thickness represents the strength of co-occurrence between two items, and color indicates cluster membership as computed by the VOS clustering algorithm.

\begin{figure}[h!]
\centering
\begin{minipage}{.80\textwidth}
\includegraphics[width=\textwidth]{images/Figure4.pdf}
\caption{Co-occurrence network of technology, DI components, mediating factors, and equity dimensions.}
\label{fig-4}
\source{VOSviewer (version 1.6.18, \textcite{vaneck2010}).}
\end{minipage}
\end{figure}

The red cluster (left) reflects classroom-level differentiation using digital tools, multimodal pro\-ducts, and inclusive learning environments under real instructional constraints. The green cluster (right) represents adaptive/AI systems, teacher professional development, and mindset, highlighting how personalization and agency depend on teacher capacity, algorithmic transparency, and infrastructural support. Bridge links (e.g., Process, Infrastructure Access) show that equitable differentiated instruction emerges from the interaction between pedagogical inclusion and adaptive technologies.

\subsection{Trends in technology-enhanced DI integration}
The review identified four key themes in the integration of technology into DI, as outlined in Table \ref{tab-3}.

%--- codigo da tabela 3 ---%
\begin{table}[h!]
\centering
\begin{threeparttable}
\small
\caption{Key themes, supporting studies, findings, and implications.}\label{tab-3}
\begin{tabular}{p{2.6cm}p{3.2cm}p{4cm}p{3cm}}
\toprule
Theme & Supporting Studies & Key Findings & Implications \\
\midrule
Theme 1: Digital Tools for Content Differentiation & \textcite{halil2024, hasanah2023, malacapay2019, yuen2023} & Teachers commonly used accessible digital tools (e.g., Google Classroom, Padlet, Quizizz, AVP, social media) to deliver multimodal content, tiered materials, and accessible resources for students with varying readiness levels. These tools enabled flexible content delivery and improved accessibility but often lacked advanced adaptivity. & Enhancing teacher digital literacy and ensuring infrastructure availability are essential for leveraging digital tools effectively to support differentiated content. \\
\addlinespace
Theme 2: Adaptive Systems and AI for Process Differentiation & \textcite{songer2019, zervoudakis2019, forster2018} & AI-powered systems such as Learning Navigator and PSO algorithms supported automated group formation, personalized learning pathways, and real-time recommendations based on learner performance. These systems reduced teacher workload but faced challenges related to transparency, data requirements, and contextual fit. & Transparency of algorithms, teacher training, and data protection protocols are required to ensure effective and ethical use of adaptive and AI-based systems in DI. \\
\addlinespace
Theme 3: Innovative Products and Inclusive Learning Environments & \textcite{wu2019, estaiteyeh2023, spyropoulou2025} & Technologies like Virtual Reality (VR), digital video games, and flexible digital products enabled students to produce diverse outputs and engage in context-rich, inclusive learning environments. Accessibility features (e.g., text-to-speech, multiple response formats) improved participation of students with diverse needs. & Teacher capacity-building for designing inclusive digital products and utilizing accessibility features is crucial for sustaining innovation and supporting diverse learners. \\
\addlinespace
Theme 4: Teacher Mindset and Contextual Factors & \textcite{halil2024, hasanah2023, yuen2023, shareefa2024} & Teacher mindset, professional development, and institutional support strongly influenced technology-enhanced DI adoption. Positive attitudes and adequate support enabled effective integration, while infrastructure limitations, rigid curricula, and time constraints hindered implementation. & School policies, teacher professional learning, and mindset shifts are critical to sustain and scale technology-supported DI practices. \\
\bottomrule
\end{tabular}
\source{Own elaboration.}
\end{threeparttable}
\end{table}

\begin{quote}
    Theme 1: Digital Tools for Content Differentiation
\end{quote}

A consistent finding across the studies \cite{halil2024, hasanah2023, malacapay2019, yuen2023} is that teachers make extensive use of accessible digital platforms such as Google Classroom, Padlet, Quizizz, and multimedia applications to differentiate learning content. These platforms support multimodal content delivery, provide tiered materials, and increase accessibility for students with varied readiness levels and diverse backgrounds. For instance, Quizizz allows teachers to conduct instant formative assessments with different difficulty levels, while Padlet facilitates collaborative curation of ideas that can be adjusted according to student ability or interest. The pedagogical reasoning behind these practices is consistent with \textcite{tomlinson1999} and \textcite{annemieke2019}, who argue that DI should begin by offering varied entry points into the content so that all learners can access the core concepts.

However, many of these tools operate at a static level of differentiation, where teachers prepare materials in advance with varying levels of complexity rather than making dynamic adjustments during the learning process. This pattern points to a gap between the potential of digital tools and how they are currently applied in classrooms. The literature emphasizes that teacher digital literacy and the availability of infrastructure are essential mediators \cite{cevikbas2022, engelbrecht2020}. Without sufficient professional development and robust digital infrastructure, the promise of differentiated content risks remaining superficial. Prior research on technology-mediated personalization \cite{shemshack2020} also cautions that digital tools must move beyond mere content dissemination to support deep conceptual understanding and learner autonomy. Thus, the challenge in this theme lies not only in equipping teachers with digital tools but also in empowering them to utilize these tools in ways that extend beyond simple content variation.

\begin{quote}
    Theme 2: Adaptive Systems and AI for Process Differentiation
\end{quote}

The integration of adaptive learning systems and AI represents the most promising yet complex frontier of DI. Studies by \textcite{songer2019, zervoudakis2019}, and \textcite{forster2018} highlight how platforms like Learning Navigator and algorithms like Particle Swarm Optimization (PSO) can automate group formation, provide personalized learning trajectories, and deliver real-time instructional recommendations. These systems embody the pedagogical principle of tailoring the learning process to each student’s performance data, closely reflecting \textcite{vygotsky1978} concept of the zone of proximal development (ZPD). By situating learners within their ZPD, AI-driven systems reduce teacher workload and allow for scalable personalization \cite{celik2022, seo2024}.

Nevertheless, the literature points to several important challenges. Transparency and algorithmic bias remain unresolved concerns \cite{chang2023, murtaza2022}. Teachers often have a limited understanding of how adaptive systems generate recommendations, which can result either in uncritical reliance or in mistrust. These systems also require large amounts of learner data, raising concerns about privacy, ethical use, and cultural relevance \cite{giannakos2025}. \textcite{heinrich2019} emphasize that successful integration depends heavily on contextual fit, since rural or under-resourced schools may face barriers such as limited bandwidth, insufficient hardware, and low levels of data literacy. The implication is that teacher training, algorithmic transparency, and data governance must accompany the adoption of AI-based DI to ensure that these tools promote equity rather than reinforce divides. This theme reveals a critical paradox. While AI and adaptive technologies have the potential to expand individualized learning at scale, their effectiveness ultimately relies on the pedagogical goals and ethical frameworks that guide their use.

\begin{quote}
    Theme 3: Innovative Products and Inclusive Learning Environments
\end{quote}

The third theme expands DI beyond content and process by focusing on student products and the creation of inclusive learning environments. Research by \textcite{wu2019, estaiteyeh2023, spyropoulou2025} demonstrates how technologies such as VR, digital video games, and flexible authoring tools enable students to produce varied and creative outputs while engaging in rich and immersive contexts. These studies align with constructivist theories of learning \cite{bruner1996, piaget1972}, which emphasize that learners deepen their understanding through the active creation and representation of knowledge in diverse forms.

Digital technologies also contribute to inclusivity by embedding accessibility features. Tools such as text-to-speech, multiple response formats, and personalized interface settings expand opportunities for participation among students with disabilities or special learning needs \cite{damyanov2024}. In this sense, DI is framed not only as a pedagogical variety but also as a mechanism of equity that allows learners with diverse needs to fully engage. This perspective is consistent with the principles of Universal Design for Learning (UDL) \cite{cast2018}, which emphasizes the importance of providing multiple means of engagement to sustain motivation and interest, multiple means of representation to present information in varied formats such as text, audio, or visuals, and multiple means of expression to allow students to demonstrate their learning in diverse ways. When DI is combined with UDL, the emphasis is placed on ensuring that technology does more than diversify instruction. It should also function as a tool to proactively remove barriers, making it possible for all learners to access and participate meaningfully in the learning process.

Yet, sustaining such innovations requires more than access to digital products. Teachers must be capacitated to design inclusive digital tasks, assess multimodal outputs, and integrate accessibility features meaningfully. Prior work \cite{molenaar2019, pozdniakov2025} stresses that teacher interpretation of digital feedback is pivotal for transforming raw technological affordances into pedagogically valuable practices. Hence, this theme illustrates that innovative products and environments cannot be decoupled from teacher expertise and the embedding of curricula.

\begin{quote}
    Theme 4: Teacher Mindset and Contextual Factors
\end{quote}
The most decisive factor in technology-supported DI is not the technology itself but the mindset and context of educators and institutions. Studies \cite{halil2024, hasanah2023, shareefa2024, yuen2023} consistently demonstrate that teachers’ attitudes toward innovation, their access to professional learning, and the level of institutional support significantly influence implementation. Teachers with a growth-oriented mindset \cite{dweck2006} are more likely to experiment with digital differentiation. In contrast, those constrained by rigid curricula, time pressures, or a lack of infrastructure tend to adopt surface-level uses.

This is consistent with earlier research \cite{smalejacobse2019, vangeel2019} showing that DI is inherently complex and depends strongly on teacher agency. Without institutional policies that provide sufficient time, resources, and incentives, even the most advanced technologies are likely to be underused. By contrast, supportive environments that encourage collaborative professional learning communities \cite{darlinghammond2020} allow teachers to share practices, build confidence, and sustain innovation. Contextual factors, therefore, play a decisive role in shaping how technology-enhanced DI can be meaningfully implemented. As \textcite{li2021} note, cultural orientations toward hierarchy, autonomy, or innovation deeply influence how teachers interpret and adopt technology-enhanced DI.

The implication is clear that scaling DI with technology requires a systemic approach that balances teacher mindset, institutional support, and cultural context. Technology alone cannot overcome structural and attitudinal barriers but must be embedded within broader educational change efforts.

\subsection{Cross-theme synthesis}
Technology-enhanced DI is more than just a set of tools; it connects pedagogy, institutional frameworks, and equity efforts. Digital tools expand access to content, but without adaptive procedures and inclusive design, they can remain static. AI and adaptive systems offer scalability, yet opaque algorithms and data-heavy processes raise equity issues. Innovative products encourage multimodality and inclusion but seldom scale without teacher expertise and policy support. Teachers’ mindsets and institutional conditions influence results across areas, highlighting that technology’s value is socially shaped and dependent on pedagogical goals and systemic backing (SDG 4).

Most evaluations emphasize short-term engagement or test improvements; long-term outcomes like autonomy, self-regulation, critical thinking, and inclusion are less studied. Algorithmic opacity and inconsistent infrastructure risk are increasing disparities. Many studies view teachers as end-users rather than partners in design. Future research should explore AI governance in education, strengthen teacher agency through ongoing professional development, and tackle issues of access, bias, and sustainability. In this view, technology acts as a strategic partner, not a replacement, promoting personalization and educational justice when aligned with equity-focused pedagogy and institutional commitment.

\section{Conclusion}
This review explored how AI, adaptive systems, and digital tools are integrated into DI to promote educational equity. Analyzing 41 studies reveals that technology-enhanced DI operates as an ecosystem connecting pedagogy, institutional frameworks, and justice goals. AI and adaptive systems increase personalized learning options, while digital platforms and inclusive products encourage multimodal engagement. However, teacher agency, ethical governance, and systemic support are crucial for making sure that technological differentiation enhances rather than damages equity.

The findings redefine DI in the digital age as a process of pedagogical mediation, where technology gains significance through teacher-guided and ethically designed implementation. Future efforts should focus on enhancing teachers’ AI literacy, promoting transparent and data-responsible practices, and incorporating digital-justice principles into institutional and policy frameworks. Long-term progress toward inclusive, technology-supported differentiation will rely on collaborative design among educators, researchers, and policymakers grounded in equity, ethics, and pedagogical purpose.

\section{Limitations and directions for future research}
While this review followed PRISMA 2020 for rigor and transparency, several limitations should be acknowledged. Limiting to English-language, peer-reviewed articles may have excluded relevant non-English or grey literature. The 2015-2025 time frame captures rapid technological changes but might omit emerging developments in generative AI. The diversity of study designs also constrained quantitative synthesis, making the findings interpretive. Future research should expand database coverage, include multilingual studies, and explore post-2025 innovations, especially regarding ethical governance, teacher AI literacy, and long-term equity outcomes.


\printbibliography\label{sec-bib}
% if the text is not in Portuguese, it might be necessary to use the code below instead to print the correct ABNT abbreviations [s.n.], [s.l.]
%\begin{portuguese}
%\printbibliography[title={Bibliography}]
%\end{portuguese}


%full list: conceptualization,datacuration,formalanalysis,funding,investigation,methodology,projadm,resources,software,supervision,validation,visualization,writing,review
\begin{contributors}[sec-contributors]
\authorcontribution{Jamil Suprihatiningrum}[conceptualization,datacuration,formalanalysis,investigation,methodology,supervision,validation,visualization,writing,review]
\authorcontribution{Ahmad Syaripudin}[datacuration,formalanalysis,investigation,projadm,resources,review]
\end{contributors}

\begin{dataavailability}
\txtdataavailability{dataonly} % options: dataavailable, dataonly, databody, datanotav, nodata
\end{dataavailability}


\end{document}

